% TEMPLATE-MILC-v3.tex - plantilla maestra para todas las guias MILC v3
% Ver scripts/generadores/build-guia-milc.py para la lista completa de
% claves esperadas. El script valida que no queden placeholders sin
% reemplazar antes de escribir el archivo final.

\documentclass[11pt,letterpaper]{article}
\usepackage[margin=1.75cm]{geometry}
\usepackage[table]{xcolor}
\usepackage{fontspec}
\usepackage[spanish]{babel}
\usepackage{tikz}
\usepackage[most]{tcolorbox}
\usepackage{tabularx}
\usepackage{array}
\usepackage{fancyhdr}
\usepackage{titlesec}
\usepackage{ragged2e}
\usepackage{enumitem}
\usepackage{microtype}

\setmainfont{Helvetica Neue}
\setsansfont{Helvetica Neue}
\setlength{\parindent}{0pt}
\setlength{\parskip}{6pt}
\hyphenpenalty=200
\exhyphenpenalty=200
\pretolerance=400
\tolerance=2000
\setlength{\emergencystretch}{1em}
\renewcommand{\arraystretch}{1.25}
\setlist{leftmargin=5mm,itemsep=1mm,topsep=1mm}
\newcolumntype{Y}{>{\RaggedRight\arraybackslash}X}

\definecolor{milcMagenta}{HTML}{E6007E}
\definecolor{milcVino}{HTML}{5A0038}
\definecolor{milcTurquesa}{HTML}{0093A5}
\definecolor{milcVerde}{HTML}{58A923}
\definecolor{milcMostaza}{HTML}{E5B400}
\definecolor{milcOcre}{HTML}{B86600}
\definecolor{milcNegro}{HTML}{241816}
\definecolor{milcGris}{HTML}{F7F7F4}
\definecolor{milcLinea}{HTML}{E8E2DD}
\definecolor{milcGradePrimary}{HTML}{0066FF}
\definecolor{milcGradeDark}{HTML}{003D99}
\definecolor{milcGradeSoft}{HTML}{D6E8FF}
\definecolor{milcGradeAccent}{HTML}{CFE7FF}
\definecolor{milcGradeText}{HTML}{FFFFFF}
\color{milcNegro}

\pagestyle{fancy}
\fancyhf{}
\lhead{\small Tecnología e Informática · Grado 11}
\rhead{\small Guía 25 · MILC v3}
\cfoot{\small\thepage}
\renewcommand{\headrulewidth}{0pt}

\newtcolorbox{titlebox}[1]{
  enhanced,colback=#1,colframe=#1,arc=7mm,boxrule=0pt,
  left=7mm,right=7mm,top=3mm,bottom=3mm,
  before skip=8pt,after skip=8pt,
  fontupper=\sffamily\bfseries\Large\color{white}
}

\newtcolorbox{softbox}[3]{
  enhanced,colback=#2,colframe=#1,borderline west={4pt}{0pt}{#1},
  arc=2mm,boxrule=.4pt,left=4mm,right=4mm,top=3mm,bottom=3mm,
  before skip=6pt,after skip=8pt,title={#3},coltitle=white,
  colbacktitle=#1,fonttitle=\sffamily\bfseries,fontupper=\RaggedRight,
  attach boxed title to top left={xshift=3mm,yshift=-2mm},
  boxed title style={arc=2mm, boxrule=0pt}
}

\newtcolorbox{drawbox}[2]{
  enhanced,colback=white,colframe=#1,arc=4mm,boxrule=1pt,
  left=5mm,right=5mm,top=4mm,bottom=4mm,before skip=8pt,after skip=8pt,
  title={#2},coltitle=white,colbacktitle=#1,fonttitle=\sffamily\bfseries,
  fontupper=\RaggedRight,
  attach boxed title to top left={xshift=4mm,yshift=-2mm},
  boxed title style={arc=3mm, boxrule=0pt}
}

\newtcolorbox{infoband}[1]{
  enhanced,colback=#1!8,colframe=#1!50,arc=2mm,boxrule=.5pt,
  left=4mm,right=4mm,top=2mm,bottom=2mm,before skip=4pt,after skip=4pt,
  fontupper=\small\RaggedRight
}

% Puente narrativo entre secciones (contrato editorial Conéctate).
% Da continuidad: "Acabas de X. Ahora vas a Y porque Z."
% Visualmente: banda izquierda en color del grado, texto en itálica pequeña.
\newtcolorbox{puentebox}{
  enhanced,colback=milcGradeSoft,colframe=milcGradeDark,
  borderline west={3pt}{0pt}{milcGradeDark},
  arc=0mm,boxrule=0pt,left=5mm,right=4mm,top=2.5mm,bottom=2.5mm,
  before skip=4pt,after skip=8pt,
  fontupper=\itshape\small\color{milcNegro}\RaggedRight
}
\newcommand{\puente}[1]{%
  \begin{puentebox}%
    \textcolor{milcGradeDark}{\textbf{\textrightarrow}}\hspace{2mm}#1%
  \end{puentebox}%
}

\newcommand{\linead}[1]{\noindent\color{milcLinea}\rule{#1}{0.5pt}\color{milcNegro}}
\newcommand{\respuesta}[1]{\par\vspace{2mm}\foreach \n in {1,...,#1}{\noindent\color{milcLinea}\rule{\linewidth}{0.5pt}\par\vspace{5mm}}\color{milcNegro}}
\newcommand{\checkbox}{\textcolor{milcGradeDark}{\fbox{\phantom{x}}}\hspace{5pt}}

\newcommand{\phasecard}[4]{
\begin{tcolorbox}[enhanced,colback=#3,colframe=#2,arc=3mm,boxrule=.5pt,height=3.05cm,valign=center,halign=center,left=1.5mm,right=1.5mm,top=2mm,bottom=2mm]
{\sffamily\bfseries\fontsize{22}{24}\selectfont\textcolor{#2}{#1}}\par
{\sffamily\bfseries\small #4}
\end{tcolorbox}
}

\begin{document}
\thispagestyle{empty}
\begin{tikzpicture}[remember picture,overlay]
  \fill[white] (current page.south west) rectangle (current page.north east);
  \path[fill=milcGradePrimary,rounded corners=16mm]
    ([xshift=.55cm,yshift=-.55cm]current page.north west) rectangle
    ([xshift=-.55cm,yshift=3.65cm]current page.south east);

  \node[anchor=north west,text=milcGradeText] at ([xshift=2.0cm,yshift=-1.85cm]current page.north west)
    {\sffamily\bfseries\fontsize{14}{16}\selectfont GRADO};
  \node[anchor=north west,text=milcGradeText,opacity=.82] at ([xshift=2.0cm,yshift=-2.75cm]current page.north west)
    {\sffamily\bfseries\fontsize{9}{11}\selectfont SERIE GUÍAS · TECNOLOGÍA E INFORMÁTICA};

  \begin{scope}[shift={([xshift=-3.05cm,yshift=-2.55cm]current page.north east)}]
    \fill[white,opacity=.80] (-.62,-.52) rectangle (.62,.52);
    \draw[milcGradeText,line width=1pt] (-.62,-.52) rectangle (.62,.52);
    \fill[milcGradeDark] (-.42,-.38) rectangle (-.22,.06);
    \fill[milcGradeAccent] (-.10,-.38) rectangle (.10,.24);
    \fill[milcGradePrimary!60!black] (.22,-.38) rectangle (.42,.38);
  \end{scope}

  \node[anchor=west,text=milcGradeText] at ([xshift=1.9cm,yshift=-12.85cm]current page.north west)
    {\sffamily\bfseries\fontsize{124}{124}\selectfont 11};
  \draw[milcGradeText,line width=1.7pt]
    ([xshift=4.52cm,yshift=-11.68cm]current page.north west) circle (.13cm);

  \node[anchor=west,text=milcGradeText,text width=8.7cm,align=left] at ([xshift=10.35cm,yshift=-11.6cm]current page.north west)
    {
     {\sffamily\bfseries\fontsize{19}{22}\selectfont Datos del MVP ---\\recoger evidencia\\con rigor}\\[4mm]
     {\sffamily\bfseries\fontsize{23}{26}\selectfont Guía 25-11-TIC}\\[2mm]
     {\sffamily\fontsize{13}{16}\selectfont Período 3 · Proyecto final emprendedor}\\[5mm]
     {\sffamily\bfseries\fontsize{11}{13}\selectfont Institución Educativa Sor María Juliana}\\[1mm]
     {\sffamily\fontsize{10}{12}\selectfont Autor: Dr. Álvaro Cárdenas Orozco}
    };

  \path[fill=milcGradeSoft,rounded corners=7mm]
    ([xshift=1.35cm,yshift=.85cm]current page.south west) rectangle
    ([xshift=-1.35cm,yshift=4.25cm]current page.south east);
  \draw[milcGradeDark,line width=.65pt,rounded corners=7mm]
    ([xshift=1.35cm,yshift=.85cm]current page.south west) rectangle
    ([xshift=-1.35cm,yshift=4.25cm]current page.south east);
  \node[anchor=center,text=milcNegro,text width=17.0cm,align=center] at ([yshift=2.55cm]current page.south)
    {\sffamily\fontsize{10}{12}\selectfont Metodología MILC · Apertura ancestral · Pensamiento computacional · Triángulo Dussel-Estoicismo-Floridi\\
    \textcolor{milcGradeDark}{\bfseries Producto:} Plan de medición del MVP (3-5 indicadores definidos como ratios o valores absolutos con meta numérica) + tabla con los primeros 10 datos reales recolectados de usuarios del MVP de la sesión 4 + análisis de 1 página con 3 hallazgos numéricos + decisión por escrito (sigo, ajusto, pivoto, descarto) anclada a los datos};
\end{tikzpicture}
\mbox{}
\newpage

\section*{Apertura: Datos del MVP --- recoger evidencia con rigor}

\begin{softbox}{milcGradeDark}{milcGradeSoft}{Saber ancestral}
En los caminos del Valle del Cauca, antes de los hospitales modernos, había una figura que dominaba un oficio que hoy llamaríamos \emph{``medicina basada en evidencia''}: \textbf{el médico tradicional, la sobandera, la curandera del barrio}. Su práctica nunca empezaba con la receta: empezaba con \textbf{observar}. La sobandera, antes de tocar la mano luxada, pedía caminar al paciente para ver la postura. La curandera, antes de dar el agua de violetas para la tos, escuchaba la respiración 3 minutos completos. El partero del campo, antes de decidir cualquier maniobra, contaba los latidos del bebé y los movimientos del vientre durante una hora. La sabiduría era rigurosa: \textbf{observar antes de actuar, medir antes de intervenir, escuchar antes de hablar}. La diferencia entre un médico tradicional respetado y un curandero novato no era el conocimiento de las plantas: era la \emph{disciplina de la observación previa}. Esa virtud ancestral es la misma que separa al emprendedor que decide con datos del que decide con corazonadas.
\end{softbox}

\begin{softbox}{milcTurquesa}{milcGris}{Saber contemporáneo}
La \textbf{medición rigurosa del MVP} es lo que separa el experimento profesional del experimento desperdiciado. Sin medición no hay aprendizaje; con medición caótica hay confusión peor que la ignorancia. La regla profesional, conocida como \textbf{ciclo Construir-Medir-Aprender} (Eric Ries), exige que cada MVP esté acompañado de 3-5 \textbf{indicadores explícitos}, definidos antes del lanzamiento, con \textbf{meta numérica} y \textbf{método de recolección} concretos. Los indicadores se distinguen en dos familias: \textbf{vanity metrics} (métricas de vanidad: número total de visitas, seguidores, likes) que se ven bien pero no dirigen decisiones, y \textbf{actionable metrics} (métricas accionables: tasa de conversión, costo de adquisición, retención semanal) que sí dirigen decisiones. La phronesis del emprendedor consiste en elegir métricas accionables sobre métricas de vanidad, aunque las segundas suenen mejor en una presentación.
\end{softbox}

\begin{softbox}{milcMostaza}{milcGris}{Pregunta-puente}
\textit{¿Qué sabía la sobandera del barrio al observar la postura del paciente durante 3 minutos antes de tocarlo, que el emprendedor novato olvida cuando salta del MVP a la \emph{``conclusión''} sin medir nada? ¿Y por qué \emph{``tuvimos 500 visitas''} es la peor manera de reportar el resultado de un MVP, aunque suene a éxito?}
\end{softbox}

\begin{softbox}{milcVerde}{milcGris}{Saber y saber hacer}
Al terminar podrás: (1) \textbf{identificar} 3-5 indicadores accionables (no de vanidad) para tu MVP, cada uno con meta numérica y método de recolección; (2) \textbf{analizar} los primeros 10 datos reales recogidos del MVP de la sesión 4 con un mínimo de rigor estadístico apropiado para muestras pequeñas; (3) \textbf{explicar} qué dicen los datos sobre el supuesto crítico que estabas probando; (4) \textbf{evaluar} la evidencia con honestidad y tomar una decisión documentada: sigo, ajusto, pivoto, descarto.
\end{softbox}

\small
\begin{tabularx}{\textwidth}{>{\bfseries}p{3.0cm}Y}
\rowcolor{milcGradePrimary!12}Autor & Dr. Álvaro Cárdenas Orozco\\
DBA & Diseño, implementación y sustentación de un proyecto de impacto local (MEN, T\&I)\\
Referentes & Dussel (analéctica · sur global) · Lean Startup · Floridi · Estoicismo\\
Producto final & Plan de medición del MVP (3-5 indicadores definidos como ratios o valores absolutos con meta numérica) + tabla con los primeros 10 datos reales recolectados de usuarios del MVP de la sesión 4 + análisis de 1 página con 3 hallazgos numéricos + decisión por escrito (sigo, ajusto, pivoto, descarto) anclada a los datos\\
\end{tabularx}
\normalsize

\puente{Antes de empezar, mira el viaje completo. La sesión 4 te dio un MVP funcional probado por al menos 1 usuario. Hoy le pones encima la \textbf{disciplina de la medición}: indicadores definidos, datos recolectados, análisis honesto. La sesión 6 tomará lo que aprendas hoy como insumo para construir el pitch. Sin datos hoy, el pitch será cuento, no evidencia.}

\begin{titlebox}{milcGradeDark}
Ruta MILC: así vamos a aprender
\end{titlebox}
\begin{tabularx}{\textwidth}{>{\centering\arraybackslash}X>{\centering\arraybackslash}X>{\centering\arraybackslash}X>{\centering\arraybackslash}X}
\phasecard{1}{milcVerde}{milcGris}{Escuta\\[-1mm]\small Observo, escucho y pregunto.} &
\phasecard{2}{milcTurquesa}{milcGris}{Sistematización\\[-1mm]\small Pensamiento computacional.} &
\phasecard{3}{milcMagenta}{milcGris}{Praxis\\[-1mm]\small Construyo y aplico.} &
\phasecard{4}{milcVino}{milcGris}{Evaluación\\[-1mm]\small Reflexión liberadora.}
\end{tabularx}


\puente{Antes de teorizar, vas a hacer una \emph{auditoría rápida} de lo que ya tienes: ¿qué datos te dejó la sesión 4?, ¿quiénes son los primeros usuarios?, ¿qué hicieron en el MVP?, ¿qué te dijeron después? Esa auditoría revela las preguntas que tu plan de medición tiene que responder.}

\begin{titlebox}{milcVerde}
Fase 1 · Escuta
\end{titlebox}

\begin{softbox}{milcVerde}{milcGris}{Escena para escuchar}
\textbf{Actividad 1 · IDENTIFICA --- 3-5 indicadores accionables para mi MVP} (20 min · individual). Toma tu MVP de la sesión 4 y define los indicadores que medirán si \emph{funcionó} respecto al supuesto crítico. Reglas: (1) cada indicador es un \textbf{ratio o un valor absoluto con meta} (no \emph{``más visitas''}, sino \emph{``20\% de los visitantes completan el formulario''} o \emph{``al menos 5 usuarios agendan demo en 7 días''}). (2) cada indicador tiene un \textbf{método de recolección} concreto (Google Analytics, conteo manual en Sheets, exportar Tally, ManyChat dashboard). (3) cubre al menos 2 de las 3 familias: \textbf{adquisición} (cuántos llegan), \textbf{activación} (cuántos hacen la acción principal), \textbf{retención} (cuántos vuelven o responden). (4) elige indicadores accionables, no de vanidad.
\end{softbox}

\checkbox Definí 3-5 indicadores con meta numérica explícita.\\
\checkbox Cada indicador tiene método de recolección concreto.\\
\checkbox Cubro al menos 2 de las 3 familias (adquisición, activación, retención).

\begin{infoband}{milcVerde}
\textbf{Lo que va al cuaderno (Actividad 1 · IDENTIFICA).} Título: «Actividad 1 --- Plan de medición de mi MVP». \textbf{Formato:} tabla de 3-5 filas × 4 columnas (Indicador~|~Familia~|~Meta numérica~|~Método de recolección). \textbf{Sabes que terminaste cuando} podrías pedirle a otro compañero que recolecte los datos siguiendo tu tabla, sin pedirte aclaraciones.
\end{infoband}

\puente{Ya tienes el material crudo. Ahora vas a entender la \textbf{anatomía del plan de medición profesional}: las 3 familias de indicadores (adquisición, activación, retención), la diferencia entre vanity y actionable metrics, los 6 errores típicos del análisis novato (muestra pequeña interpretada como definitiva, sesgo de confirmación, datos sin baseline, indicador único, métricas de vanidad como prueba).}

\begin{titlebox}{milcTurquesa}
Fase 2 · Sistematización con pensamiento computacional
\end{titlebox}

\begin{softbox}{milcTurquesa}{milcGris}{Cuatro pilares aplicados al tema}
Un plan de medición sin teoría detrás es lista de números sueltos. Para que tu medición tenga rigor profesional necesitas la \textbf{anatomía del marco AAR} (Adquisición, Activación, Retención) y la \textbf{distinción crítica} entre métricas de vanidad y métricas accionables. Estas dos piezas conceptuales separan al analista del coleccionista de números.
\end{softbox}

\small
\begin{tabularx}{\textwidth}{>{\bfseries\RaggedRight\arraybackslash}p{3.6cm}Y}
\rowcolor{milcTurquesa!90}\color{white}Pilar & \color{white}\bfseries Aplicación al tema\\
1. Descomposición & Las métricas se organizan en 3 familias del modelo \textbf{AARRR} (popularizado por Dave McClure y simplificado aquí a las 3 más críticas para MVP): (1) \textbf{Adquisición}: cuántas personas llegan al MVP (visitas únicas, clicks en link, mensajes recibidos). Responde: ¿el canal funciona? (2) \textbf{Activación}: cuántas hacen la acción principal del MVP (completar formulario, agendar demo, descargar, conversar con bot). Responde: ¿la propuesta de valor enganchó? (3) \textbf{Retención}: cuántas vuelven o responden cuando las contactas después (apertura de correo de seguimiento, respuesta a WhatsApp, reuso del MVP). Responde: ¿hubo valor real o curiosidad pasajera?\\
2. Reconocimiento de patrones & La \textbf{distinción vanity vs actionable} es la línea que separa al profesional del aficionado: (a) \textbf{Vanity metric}: número grande que se ve bien pero no permite decidir nada (\emph{``500 visitas''}, \emph{``100 likes''}, \emph{``mucha gente vio mi post''}). (b) \textbf{Actionable metric}: ratio o valor que dirige una decisión concreta (\emph{``de 100 visitantes, 12 completaron formulario = 12\%''}, \emph{``costo por usuario adquirido = \$3.000''}, \emph{``retención semana 1 = 40\%''}). La phronesis profesional consiste en \textbf{reportar siempre las accionables}, aunque las de vanidad suenen mejor.\\
3. Abstracción & Lo esencial cabe en una pregunta: \emph{¿esta métrica me hace cambiar lo que haría mañana?}. Si la respuesta es no, es métrica de vanidad. Si la respuesta es sí (\emph{``si la activación es menor a 10\%, cambio el copy''}, \emph{``si la retención semana 1 es menor a 30\%, repienso el valor''}), es accionable. Esa simple pregunta filtra el 80\% de las métricas que circulan en presentaciones de emprendedores.\\
4. Algoritmo conceptual & Algoritmo de análisis para 10 datos (muestra pequeña): (1) recoge los datos en una tabla (1 fila por usuario, columnas con los indicadores); (2) calcula los ratios de cada familia (no solo totales); (3) marca con asterisco los usuarios que se comportaron como esperabas; (4) busca patrones cualitativos (¿hay algo común en los que sí, en los que no?); (5) acepta la regla estadística básica: con n=10 no se hacen afirmaciones definitivas, se identifican \emph{tendencias} que sugieren la siguiente prueba; (6) declara explícitamente las limitaciones de la muestra en tu informe.\\
\end{tabularx}
\normalsize

\begin{drawbox}{milcTurquesa}{Anatomía del plan de medición}
\textbf{Adquisición:} cuántos llegan + meta. \\[1mm] \textbf{Activación:} cuántos hacen la acción principal + meta. \\[1mm] \textbf{Retención:} cuántos vuelven o responden + meta. \\[1mm] \textbf{Método:} dónde se recoge cada dato. \\[1mm] \textbf{Baseline:} valor inicial de referencia. \\[1mm] \textbf{Umbral de decisión:} qué resultado disparará qué acción.
\end{drawbox}

\begin{softbox}{milcMostaza}{milcGris}{Errores comunes y su corrección}
(1) \textbf{Muestra pequeña interpretada como definitiva}: con n=10 no se concluye, se sugiere. (2) \textbf{Sesgo de confirmación}: leer los datos buscando la conclusión que ya quieres. (3) \textbf{Datos sin baseline}: reportar \emph{``20 inscritos''} sin decir si eso es mucho o poco respecto a la meta. (4) \textbf{Indicador único}: medir solo \emph{``visitas''} y no medir activación ni retención (te quedas en adquisición). (5) \textbf{Métricas de vanidad como prueba}: usar likes o seguidores como evidencia de validación. (6) \textbf{Mezclar datos de distintos canales} sin separar: si 5 vienen de Instagram y 5 de WhatsApp, son dos pruebas distintas. (7) \textbf{No documentar el contexto}: no anotar qué día se lanzó, qué cambió en el camino, qué pasó en el entorno (la fecha y el contexto cambian todo).
\end{softbox}

\begin{infoband}{milcTurquesa}
\textbf{Lo que va al cuaderno (Actividad 2 · ANALIZA + EXPLICA).} Título: «Actividad 2 --- Anatomía de la medición». \textbf{Formato:} (a) los 6 campos de la anatomía aplicados a tu MVP; (b) lista de los 7 errores con marca del que más temes; (c) tabla preparada para llenar con 10 filas y columnas de tus indicadores. \textbf{Sabes que terminaste cuando} la tabla está lista para recibir los datos esta semana.
\end{infoband}

\puente{Tienes el método. Toca aplicarlo: redactar el plan de medición con 3-5 indicadores, recolectar los primeros 10 datos reales (durante la semana), procesarlos en una tabla, escribir 3 hallazgos numéricos y tomar la decisión.}

\begin{titlebox}{milcMagenta}
Fase 3 · Praxis con pensamiento computacional
\end{titlebox}

\begin{softbox}{milcMagenta}{milcGris}{Construyendo el producto con los cuatro pilares}
\textbf{Actividad 3 · EVALÚA --- Análisis y decisión basada en datos} (30 min en sesión, recolección a lo largo de la semana · individual). Esta actividad se completa al final de la semana, cuando los 10 primeros datos están en la tabla. Aquí no se inventa: se procesa lo recogido, se nombran 3 hallazgos numéricos y se toma una decisión documentada respaldada por la evidencia.
\end{softbox}

\small
\begin{tabularx}{\textwidth}{>{\bfseries\RaggedRight\arraybackslash}p{3.6cm}Y}
\rowcolor{milcMagenta!90}\color{white}Pilar & \color{white}\bfseries Aplicación al producto\\
Descomposición & Tu análisis se estructura en 4 bloques de informe: (1) \textbf{Datos crudos}: la tabla con los 10 usuarios y sus comportamientos en cada indicador. (2) \textbf{Cálculos}: ratios y medias por indicador (adquisición \%, activación \%, retención \%). (3) \textbf{3 hallazgos numéricos}: las 3 frases más importantes que dicen los datos (ejemplo: \emph{``de 10 visitantes, 4 completaron formulario = 40\% activación''}, \emph{``de 10 usuarios, 3 respondieron al correo de seguimiento = 30\% retención''}, \emph{``80\% de los inscritos vinieron de WhatsApp, 20\% de Instagram''}). (4) \textbf{Decisión}: \emph{sigo} (el MVP confirma el supuesto), \emph{ajusto} (la dirección es correcta, hay que mejorar X), \emph{pivoto} (los datos sugieren cambiar el supuesto), \emph{descarto} (la evidencia es contraria).\\
Patrones & Sigue el patrón \textbf{``hallazgo sobre interpretación''}: en el análisis, prioriza la frase del hallazgo numérico sobre la interpretación. La interpretación viene después y es honestamente provisional con n=10. Por ejemplo: \emph{``Hallazgo: 4/10 completaron el formulario. Interpretación posible: la activación está en un rango razonable para MVP en etapa temprana; conviene confirmar con n=30 antes de decidir''}.\\
Abstracción & Las 4 decisiones tienen criterios profesionales claros: (a) \textbf{Sigo} si los indicadores cumplieron al menos 70\% de las metas. (b) \textbf{Ajusto} si los indicadores están cerca pero falta refinamiento (50-70\% de meta). (c) \textbf{Pivoto} si los indicadores fallaron pero los datos revelan otro patrón (alguien se inscribió pero por razón distinta a la esperada, algún segmento no previsto respondió). (d) \textbf{Descarto} si los indicadores fallaron y los datos no revelan caminos alternativos. \textbf{Pivotar no es fracasar}; descartar a tiempo tampoco lo es. El fracaso real es seguir adelante sin evidencia.\\
Algoritmo de armado & Algoritmo: (1) calcula los ratios; (2) compara con las metas; (3) escribe los 3 hallazgos numéricos en orden de relevancia; (4) declara las limitaciones (n pequeño, sesgo de canal, sesgo de selección); (5) toma la decisión con justificación de 2 líneas anclada en números, no en sensaciones; (6) si pivotas o ajustas, formula el nuevo supuesto y el plan para la siguiente iteración.\\
\end{tabularx}
\normalsize

\begin{drawbox}{milcMagenta}{Checklist antes de cerrar el análisis}
\checkbox 10 datos reales recogidos durante la semana. \\[1mm] \checkbox Tabla con 1 fila por usuario y columnas de indicadores. \\[1mm] \checkbox Ratios calculados (no solo totales). \\[1mm] \checkbox 3 hallazgos numéricos escritos como frases concretas. \\[1mm] \checkbox Limitaciones de la muestra declaradas explícitamente. \\[1mm] \checkbox Decisión explícita (sigo, ajusto, pivoto, descarto) con justificación numérica. \\[1mm] \checkbox Informe cabe en 1 página.
\end{drawbox}

\begin{softbox}{milcMagenta}{milcGris}{Plantilla de trabajo}
\textbf{Datos crudos.} \textit{Tabla 10 filas × N columnas.} \\[1mm] \textbf{Cálculos.} \textit{Adquisición \% + Activación \% + Retención \%.} \\[1mm] \textbf{Hallazgo 1.} \textit{Frase numérica relevante.} \\[1mm] \textbf{Hallazgo 2.} \textit{Frase numérica relevante.} \\[1mm] \textbf{Hallazgo 3.} \textit{Frase numérica relevante.} \\[1mm] \textbf{Limitaciones.} \textit{Tamaño de muestra y sesgos.} \\[1mm] \textbf{Decisión.} \textit{Sigo / Ajusto / Pivoto / Descarto + justificación numérica.}
\end{softbox}

\begin{infoband}{milcMagenta}
\textbf{Lo que va al cuaderno (Actividad 3 · EVALÚA).} Título: «Actividad 3 --- Análisis del MVP». \textbf{Formato:} 1 página con tabla de 10 filas + cálculos + 3 hallazgos + limitaciones + decisión documentada. \textbf{Sabes que terminaste cuando} podrías mostrarle el informe a un consultor de datos y defender la decisión solo con la tabla en mano.
\end{infoband}

\puente{Lo que sigue resume \emph{qué} entregas al final de la semana: plan de medición + tabla de 10 datos reales + análisis de 1 página + decisión documentada. El criterio principal: que un analista de datos junior, mirando tu tabla, pueda confirmar o cuestionar tu decisión sin tener que confiar en tu palabra.}

\begin{titlebox}{milcMostaza}
Producto de la sesión
\end{titlebox}
\textbf{Plan de medición + 10 datos reales + análisis con 3 hallazgos + decisión documentada}.
\begin{softbox}{milcMostaza}{milcGris}{Criterios de calidad}
(1) 3-5 indicadores accionables con meta y método. (2) 10 datos reales recolectados de usuarios del MVP. (3) Ratios calculados por indicador. (4) 3 hallazgos numéricos escritos como frases concretas. (5) Limitaciones de la muestra declaradas explícitamente. (6) Decisión documentada (sigo, ajusto, pivoto, descarto) con justificación numérica.
\end{softbox}

\puente{Antes de cerrar, mira la medición desde las cinco dimensiones humanas. Recoger evidencia con rigor es respeto por el oficio entero; saltar a conclusiones sin datos es engaño al docente, al usuario y a uno mismo.}

\begin{titlebox}{milcVino}
Fase 4 · Evaluación liberadora — cinco dimensiones
\end{titlebox}

\begin{softbox}{milcVino}{milcGris}{Sentido de la evaluación}
Evaluar no es solo revisar una nota. Es reconocer qué aprendí, qué cambió en mí, qué emociones manejé, cómo me relaciono con la ciudad y la comunidad, y qué le dejaré a quien viene detrás.
\end{softbox}

\textbf{Mi reflexión liberadora en cinco dimensiones.} Completa cada frase iniciada:

\small
\begin{tabularx}{\textwidth}{>{\bfseries\RaggedRight\arraybackslash}p{3.4cm}Y>{\RaggedRight\arraybackslash}p{4.6cm}}
\rowcolor{milcVino}\color{white}Dimensión & \color{white}\bfseries Mi respuesta (frame starter) & \color{white}\bfseries Algo que descubrí\\
1. Personal & Lo que más cambió en mí fue \linead{6.5cm} & \linead{4.2cm}\\
2. Emocional & La emoción más fuerte fue \linead{2.5cm} y la regulé \linead{3cm} & \linead{4.2cm}\\
3. Ciudadana & En democracia esto importa porque \linead{6cm} & \linead{4.2cm}\\
4. Local (Cartago) & En mi barrio sería útil para \linead{6.5cm} & \linead{4.2cm}\\
5. Intergeneracional & A mi abuela/abuelo le diría \linead{6.5cm} & \linead{4.2cm}\\
\end{tabularx}
\normalsize

\begin{infoband}{milcVino}
\textbf{Para la próxima sustentación.} Vuelve a esta tabla en seis meses. Verás que las respuestas cambian — no porque lo aprendido cambie, sino porque tú cambias. La evaluación liberadora es una conversación contigo a través del tiempo.
\end{infoband}


\puente{Y para terminar, tres voces sobre el rigor y la verdad. Dussel sobre las decisiones que respetan a quien sufre el problema, Marco Aurelio sobre la disciplina de ver las cosas como son, Floridi sobre la ética del dato bien recolectado en la era de la información.}

\begin{titlebox}{milcGradeDark}
Triángulo de pensamiento · cierre filosófico
\end{titlebox}

\begin{softbox}{milcVino}{milcGris}{Dussel · lente del nosotros}
\textit{``Decidir sobre el otro sin escuchar el dato que el otro produce es repetir la lógica del poder sin verdad.''}

Tu MVP recogió datos de usuarios reales. Esos datos no son tuyos: son la voz numérica de personas que te confiaron 2 minutos de su tiempo. La filosofía de la liberación pide tomarlos en serio aunque contradigan tu hipótesis. Cuando un emprendedor descarta datos adversos porque \emph{``la muestra era pequeña''} mientras acepta datos favorables sin objeción, repite la lógica colonial: solo escucha al otro cuando le da la razón. La phronesis del oficio consiste en \textbf{aceptar el dato con el mismo rigor cuando confirma que cuando contradice}.

\textbf{¿Estoy tratando los datos adversos con el mismo rigor que los favorables, o solo acepto la evidencia que me da la razón?}
\end{softbox}
\linead{\linewidth}\\[1mm]
\linead{\linewidth}

\begin{softbox}{milcMostaza}{milcGris}{Estoicismo · Marco Aurelio · cuidado interior}
\textit{``Ve las cosas como son, no como te gustaría que fueran; ese es el principio del juicio recto.''}

Es tentador inflar los hallazgos numéricos para que el MVP parezca exitoso, o subestimarlos para tener excusa de pivotar. La disciplina estoica recomienda lo contrario: \emph{describe los números como son, ni más grandes ni más pequeños}. Esa precisión es virtud profesional. El emprendedor que se engaña con sus propios datos toma decisiones sobre arena. La sesión 6 (pitch) y las siguientes castigan sin piedad esa autoindulgencia con un fracaso público.

\textbf{¿Estoy interpretando los datos como son, o como me conviene que sean para que mi MVP parezca exitoso ante el docente?}
\end{softbox}
\linead{\linewidth}\\[1mm]
\linead{\linewidth}

\begin{softbox}{milcTurquesa}{milcGris}{Floridi · lente de la infoesfera}
\textit{``Recolectar datos con rigor y limitaciones declaradas es la nueva ética del oficio en la era de la información.''}

La era de la información produce datos por toneladas, pero la calidad ética del dato depende de cómo se recolecta y cómo se comunica. Un análisis profesional declara siempre: el tamaño de la muestra, el método de recolección, los sesgos posibles, las limitaciones. Esa transparencia es virtud técnica contemporánea. Los análisis opacos producen decisiones malas con apariencia de certeza, que es la peor combinación posible.

\textbf{¿Mi informe declara honestamente las limitaciones (n=10, sesgo de canal, sesgo de selección) o las esconde para sonar más convincente?}
\end{softbox}
\linead{\linewidth}\\[1mm]
\linead{\linewidth}

\begin{titlebox}{milcVino}
Compromiso final
\end{titlebox}

\small
\begin{tabularx}{\textwidth}{>{\RaggedRight\arraybackslash}p{3.0cm}YYY}
\rowcolor{milcVino}\color{white}\bfseries Criterio & \color{white}\bfseries Lo logré & \color{white}\bfseries En proceso & \color{white}\bfseries Necesito apoyo\\
Comprensión del tema & Explico ideas centrales con mis palabras. & Reconozco ideas pero falta precisión. & Necesito repasar conceptos base.\\
Pensamiento computacional & Apliqué los 4 pilares al diseñar mi producto. & Apliqué algunos pilares. & Necesito releer la fase 2.\\
Aplicación y producto & Mi producto cumple los criterios y se puede revisar. & Producto funcional con ajustes pendientes. & Aún en construcción.\\
Reflexión 5 dimensiones & Respondí cada dimensión con honestidad. & Respondí algunas con detalle. & Mis respuestas son muy generales.\\
Triángulo de pensamiento & Conecté las 3 voces con mi trabajo real. & Conecté al menos una. & No relaciono las voces con mi proceso.\\
\end{tabularx}
\normalsize

\textbf{Hoy entendí que:}\\[1mm]
\linead{\linewidth}\\[1mm]
\linead{\linewidth}

\textbf{Mi compromiso para la próxima sesión es:}\\[1mm]
\linead{\linewidth}\\[1mm]
\linead{\linewidth}

\begin{softbox}{milcMostaza}{milcGris}{Cita de despedida · Marco Aurelio}
\textit{``El obstáculo en el camino se vuelve el camino. Lo que estorba al camino se vuelve camino.''}\\[1mm]
Cada error de hoy, bien observado, se convierte en pieza del camino que pisarás mañana.
\end{softbox}

\small
\begin{tabularx}{\textwidth}{>{\bfseries}p{3.6cm}Y}
\rowcolor{milcGradeDark!12}Nombre completo: & \linead{10cm}\\
Fecha: & \linead{4cm} \quad Curso/grupo: \linead{4cm}\\
Mi firma: & \linead{9cm}\\
\end{tabularx}
\normalsize

\begin{infoband}{milcGradeDark}
\textbf{Para la próxima sesión.} Llega con el plan de medición, los 10 datos y la decisión documentada. En la sesión 6 vas a construir el \textbf{pitch}: 7 slides + 3 minutos hablados que cuenten la historia del proyecto. Los hallazgos de hoy son la columna vertebral del pitch. Sin datos hoy, el pitch será cuento; con datos hoy, será evidencia.
\end{infoband}

\end{document}
